% ============================================================
% f2014_shaanxi_arts_q10_cubic.tex
% 2014 陕西卷（文科）第 10 题：三次函数弯曲路段与两条直道相切
% 来源 PDF：普通高考/2014/2014陕西文.pdf
% 原页：第 1 页（A3），印面 913，第 10 题题干下方
% 弯曲路段为 y=x^3/2-x^2/2-x（0≤x≤2，虚线），
% 在 O(0,0) 与 y=-x 相切，在 (2,0) 与 y=3x-6 相切；
% 灰色区域为湖面，其下沿为不规则岸线，岸线不是函数图象。
% ============================================================
\begin{tikzpicture}[>=Stealth, line width=0.7pt, font=\small,
  x=1.05cm, y=0.82cm, line cap=round, line join=round]
  % 湖面（灰色区域，下沿为不规则岸线）
  \fill[gray!22]
    (-2.55,3.00) -- (2.95,3.00) -- (2.75,2.70) -- (2.55,2.40)
    -- (2.40,1.85) -- (2.15,1.75) -- (2.00,1.35) -- (1.70,1.25)
    -- (1.45,1.00) -- (1.10,0.95) -- (0.75,0.78) -- (0.40,0.72)
    -- (0.05,0.90) -- (-0.30,1.05) -- (-0.60,1.25) -- (-0.85,1.55)
    -- (-1.10,1.65) -- (-1.30,1.95) -- (-1.55,2.05) -- (-1.75,2.45)
    -- (-1.95,2.55) -- cycle;
  \draw[line width=0.6pt]
    (2.95,3.00) -- (2.75,2.70) -- (2.55,2.40)
    -- (2.40,1.85) -- (2.15,1.75) -- (2.00,1.35) -- (1.70,1.25)
    -- (1.45,1.00) -- (1.10,0.95) -- (0.75,0.78) -- (0.40,0.72)
    -- (0.05,0.90) -- (-0.30,1.05) -- (-0.60,1.25) -- (-0.85,1.55)
    -- (-1.10,1.65) -- (-1.30,1.95) -- (-1.55,2.05) -- (-1.75,2.45)
    -- (-1.95,2.55) -- (-2.55,3.00);
  \node at (0.55,2.05) {湖面};

  % 坐标轴
  \draw[->] (-4.40,0) -- (4.20,0) node[right] {$x$（千米）};
  \draw[->] (0,-2.65) -- (0,3.35) node[above] {$y$（千米）};
  \node[below left] at (0,0) {$O$};

  % 两条直道
  \draw (-3.40,3.40) -- (0,0);
  \draw (2,0) -- (3.30,3.90);
  \node[left] at (-3.15,2.60) {$y=-x$};
  \node[right] at (3.02,3.02) {$y=3x-6$};

  % 弯曲路段（三次函数图象）
  \draw[dashed, thick, domain=0:2, samples=100]
    plot (\x,{0.5*(\x)^3-0.5*(\x)^2-(\x)});

  % 切点与横坐标标注
  \fill (0,0) circle (1.2pt);
  \fill (2,0) circle (1.2pt);
  \node[below right, inner sep=1pt] at (2,0) {$2$};
\end{tikzpicture}
